\documentclass[11pt]{article}
\usepackage[margin=1in]{geometry}
\usepackage{amsmath}
\usepackage{amssymb}
\usepackage{enumitem}
\usepackage[hidelinks]{hyperref}
\usepackage{xcolor}

\title{CSCE 355: JFLAP4 --- NFAs and $\lambda$-NFAs}
\author{}
\date{}

\renewcommand{\labelitemi}{$\bullet$}
\newcommand{\fn}[1]{\texttt{\detokenize{#1}}}
\newcommand{\re}[1]{\texttt{#1}}

\begin{document}
\maketitle

\section*{Overview}

In this assignment you will build seven nondeterministic finite automata in JFLAP: five NFAs
and two $\lambda$-NFAs (NFAs with empty-string transitions, also called $\varepsilon$-NFAs).
Each is checked automatically. The assignment is worth \textbf{30 points}.

Regular expressions below use course notation: \re{+} is union, \re{*} is Kleene star, and
$\lambda$ is the empty string.

\subsection*{Getting started}

\begin{itemize}
  \item Read the JFLAP tutorial: \url{https://jflap.org/tutorial/}.
  \item JFLAP is a free \texttt{.jar} file from \url{https://www.jflap.org}. You need
        Java 8 or higher. Start it with \texttt{java -jar JFLAP7\_1.jar} (use the exact name of
        the file you downloaded; on macOS and Linux, upper and lower case matter).
  \item Unzip this assignment folder. (It is a good idea to append your username to the
        folder name, e.g.\ \fn{JFLAP4_jsmith}.)
  \item In JFLAP use \textbf{File $\to$ Open} to open one of the provided \texttt{.jff}
        files. Each one is an empty \emph{Finite Automaton}. Build your machine, then
        \textbf{File $\to$ Save} it \emph{under the same file name}. Do not change the
        machine type.
  \item JFLAP shows the empty string as $\lambda$. Just leave it as $\lambda$; do not change it
        to $\varepsilon$. To make a $\lambda$ transition, draw the transition and press Enter
        without typing anything.
\end{itemize}

\begin{center}
\begin{tabular}{|c|l|l|c|}
\hline
\textbf{\#} & \textbf{Your machine (edit this file)} & \textbf{Test strings (do not edit)} & \textbf{Points}\\
\hline
1 & \fn{problem1_01then10.jff}          & \fn{problem1_01then10_tests.txt}          & 3 \\
2 & \fn{problem2_cab.jff}               & \fn{problem2_cab_tests.txt}               & 4.5 \\
3 & \fn{problem3_no010.jff}             & \fn{problem3_no010_tests.txt}             & 4.5 \\
4 & \fn{problem4_last_new.jff}          & \fn{problem4_last_new_tests.txt}          & 3 \\
5 & \fn{problem5_10i1.jff}              & \fn{problem5_10i1_tests.txt}              & 3 \\
6 & \fn{problem6_nfa_01_110_011.jff}    & \fn{problem6_nfa_01_110_011_tests.txt}    & 4.5 \\
7 & \fn{problem7_enfa_01_110_011.jff}   & \fn{problem7_enfa_01_110_011_tests.txt}   & 7.5 \\
\hline
\end{tabular}
\end{center}

\section*{NFA or $\lambda$-NFA? Rules for every machine}

\noindent\begin{minipage}{\linewidth}
\setlength{\parindent}{0pt}
The test strings cannot tell an NFA from a $\lambda$-NFA, because both kinds of machine
recognize the same languages. What separates them here is the machine itself, and that is
checked automatically, before any test string is run:

\begin{center}
\begin{tabular}{|c|l|l|l|}
\hline
\textbf{Problem} & \textbf{Build a} & \textbf{$\lambda$ transitions} & \textbf{Size limit}\\
\hline
1 & $\lambda$-NFA & \textbf{required}, and must matter & at most 4 states \\
2 & NFA & \textbf{not allowed} & at most 4 states \\
3 & NFA & \textbf{not allowed} & at most 6 states \\
4 & NFA & \textbf{not allowed} & none \\
5 & NFA & \textbf{not allowed} & at most 8 states \\
6 & NFA & \textbf{not allowed} & at most 8 states \\
7 & $\lambda$-NFA & \textbf{required}, and must matter & at most 12 edges \\
\hline
\end{tabular}
\end{center}
\end{minipage}

A machine that breaks one of these rules receives \textbf{zero credit} for that problem, even if
it accepts and rejects every string correctly:
\begin{itemize}
  \item \textbf{NFA problems (2--6): no $\lambda$ transitions at all.} A single $\lambda$
        transition anywhere in the file earns 0 for that problem.
  \item \textbf{$\lambda$-NFA problems (1 and 7): use $\lambda$ transitions, and make them
        matter.} The grader deletes all of your $\lambda$ transitions and checks that the
        language changes. A correct machine with no $\lambda$ transitions, or with $\lambda$
        transitions that could simply be removed, earns 0 for that problem.
  \item \textbf{One symbol per transition.} Every other transition reads exactly one symbol of
        the problem's alphabet. A label such as \texttt{01} (two symbols in a row) earns 0.
        Draw a separate transition for each symbol: a self-loop on \texttt{a}, \texttt{b} and
        \texttt{c} is three transitions, and JFLAP stacks their labels above the state.
        \emph{Comma labels:} the grader reads a label such as \texttt{0,1} as one transition
        per symbol, and each one counts separately toward a size limit (\texttt{0,1} is two
        edges, \texttt{a,b,c} is three). But JFLAP itself does not: its Multiple Run reads
        \texttt{0,1} as a single three-character string, so your JFLAP tests will not match.
        Use separate transitions.
  \item \textbf{Exactly one initial state.}
\end{itemize}
Nondeterminism is of course allowed in both kinds of machine: a state may have several
transitions on the same symbol, or none at all (a string that gets stuck is rejected, so you do
not need a trap state).

\textbf{Size limits.} \emph{Every} state in your file counts toward a state limit, and
\emph{every} transition (including self-loops and $\lambda$ transitions, with a comma label
such as \texttt{0,1} counting once per symbol) counts toward Problem 7's edge limit, whether or not it can be reached, so delete anything you are not using. Each
state or edge over the limit costs \textbf{50\%} of that problem's points: one over the limit
earns half credit, two over earns nothing. Every limit can be met: there is a correct machine
within each one.

\section*{Problems}

\subsection*{Problem 1: $\lambda$-NFA (3 points, at most 4 states)}
Construct a $\lambda$-NFA, in no more than 4 states, that recognizes any number of \texttt{01}s
followed by any number of \texttt{10}s, without overlap: the language \re{(01)*(10)*}. For
example, \texttt{010} is \emph{not} in the language. The alphabet is $\{0,1\}$.
\begin{itemize}
  \item Your machine must use $\lambda$ transitions productively: deleting all of them must
        change the language it recognizes.
  \item The test file has 8 strings: the first 3 should accept, the last 5 should reject.
  \item The empty string should accept (see ``Testing the empty string'' below).
\end{itemize}

\subsection*{Problem 2: NFA (4.5 points, at most 4 states)}
Construct an NFA (no $\lambda$ transitions) in 4 or fewer states that recognizes the strings that contain \texttt{cab} as a
substring: \re{(a+b+c)*cab(a+b+c)*}. The alphabet is $\{a,b,c\}$.
\begin{itemize}
  \item The test file has 8 strings: the first 4 should accept, the last 4 should reject.
  \item The empty string should reject.
  \item Remember that a self-loop on \texttt{a}, \texttt{b} and \texttt{c} is three separate
        transitions, one per symbol.
\end{itemize}

\subsection*{Problem 3: NFA (4.5 points, at most 6 states)}
Construct an NFA (no $\lambda$ transitions) in 6 or fewer states that recognizes the strings that do not contain
\texttt{010} as a substring, \emph{or} that are just \texttt{01} repeated, that is, in
\re{(01)*}. For example, \texttt{01010101} is accepted but \texttt{0101010} is not. The empty
string should be accepted. The alphabet is $\{0,1\}$.
\begin{itemize}
  \item The test file has 8 strings: the first 4 should accept, the last 4 should reject.
  \item The empty string should accept.
\end{itemize}

\subsection*{Problem 4: NFA (3 points, no size limit)}
Construct an NFA (no $\lambda$ transitions) that recognizes the strings over $\Sigma=\{1,2,3,4\}$ whose last symbol does
not appear anywhere before it. For example, \texttt{1234} and \texttt{4231} are accepted;
\texttt{12344} and \texttt{4214} are not.
\begin{itemize}
  \item You may accept the empty string (it will not be tested).
  \item The test file has 10 strings: the first 5 should accept, the last 5 should reject.
  \item Remember to add a separate transition for each input symbol.
\end{itemize}

\subsection*{Problem 5: NFA (3 points, at most 8 states)}
Construct an NFA (no $\lambda$ transitions) in 8 or fewer states that recognizes the strings over $\Sigma=\{0,1,2\}$ that
contain a substring $1\,0^i\,1$ (a \texttt{1}, then $i$ copies of \texttt{0}, then a
\texttt{1}) where $i > 0$ is a multiple of 2 or a multiple of 3, that is,
$i \in \{2, 3, 4, 6, 8, 9, 10, 12, \ldots\}$:
\[
  \re{(0+1+2)*\,1\,(00(00)* + 000(000)*)\,1\,(0+1+2)*}
\]
For example, \texttt{1001} ($i=2$), \texttt{10001} ($i=3$), \texttt{100001} ($i=4$) and
\texttt{21210012121} are accepted; \texttt{101} ($i=1$), \texttt{1000001} ($i=5$) and
\texttt{2100000001} ($i=7$) are not.
\begin{itemize}
  \item The test file has 17 strings: the first 9 should accept, the last 8 should reject.
  \item The empty string should reject.
\end{itemize}

\subsection*{Problem 6: NFA (4.5 points, at most 8 states)}
Construct an NFA, with no $\lambda$ transitions and in no more than 8 states, recognizing the
concatenation of any number of \texttt{01}, \texttt{110}, and/or \texttt{011}, zero or more
copies of each in any order (the copies may be interleaved; look at the tests):
\re{(01+110+011)*}. The alphabet is $\{0,1\}$.
\begin{itemize}
  \item The test file has 117 strings: the first 107 should accept, the last 10 should reject.
  \item The empty string should accept.
\end{itemize}

\subsection*{Problem 7: $\lambda$-NFA (7.5 points, at most 12 edges)}
Construct a $\lambda$-NFA for the same language as Problem 6, \re{(01+110+011)*}. This is easy,
actually, but this time the limit is on \emph{edges}, not states: no more than 12 edges. Every
transition counts as one edge, including self-loops and $\lambda$ transitions; a label such as
\texttt{0,1} counts as two edges.
\begin{itemize}
  \item Your machine must use $\lambda$ transitions productively: deleting all of them must
        change the language it recognizes. (A $\lambda$ transition that could simply be removed
        does not count.)
  \item The test file is the same as Problem 6's: the first 107 strings should accept, the last
        10 should reject.
  \item The empty string should accept.
  \item Not graded: compare your solutions to Problems 6 and 7 with respect to number of states,
        number of edges, clarity, etc.
\end{itemize}

\section*{Testing your machines}

\subsection*{In JFLAP}
\begin{enumerate}
  \item Open your machine (it needs an initial state to be tested).
  \item Choose \textbf{Input $\to$ Multiple Run}.
  \item Click \textbf{Load Inputs} (at the bottom) and select the problem's test file.
  \item Click \textbf{Run Inputs} and compare each row with what the problem says should
        accept and reject.
\end{enumerate}

\subsubsection*{Testing the empty string}
The test files do not contain the empty string. In Multiple Run, put the cursor on the last
row and press Enter to add a blank line (an empty input), then run again. The empty string
should accept in Problems 1, 3, 6 and 7, and reject in Problems 2 and 5. (Problem 4 does not
test it.)

\subsection*{From the command line}
This folder contains \texttt{test\_nfa.py}, the same script that is used to grade you (Python 3
is required). From inside this folder run
\begin{center}
\texttt{python3 test\_nfa.py}
\end{center}
to check all seven problems, or \texttt{python3 test\_nfa.py 3} to check only Problem 3. The
report tells you, for each problem: whether the machine is the required kind and follows the rules above, which provided
test strings are wrong, its size, and your score. Run it as often as you like.

\section*{Grading}

Each problem earns its points only if \emph{all} of the following hold:
\begin{itemize}
  \item The machine is the required kind (NFA or $\lambda$-NFA) and follows the rules above.
  \item Every string in the problem's test file gets the expected result, and so does the
        empty string where the problem says what it should do.
  \item Your machine recognizes \emph{exactly} the described language. The provided strings are
        only examples: the grader compares your machine with the described language on all
        strings, so a machine that passes the samples but is wrong somewhere else earns no
        credit. When this happens the report shows a shortest string on which your machine is
        wrong.
\end{itemize}
A machine that meets all of these but is over its size limit loses 50\% of the problem's points
per extra state (Problem 7: per extra edge).

\textbf{Do not edit the test files or \texttt{test\_nfa.py}.} They are replaced with the
original versions before grading, so changing them can only make your own testing misleading.

\section*{Submission}

\begin{enumerate}
  \item Open \textbf{YOURUSERNAME.txt} and put your university username as the \emph{first
        line} of the file (only your username, nothing else on that line). This is how your
        submission is identified.
  \item Make sure each of your seven machines is saved in its provided \texttt{.jff} file. Do
        not rename any of the provided files.
  \item Put screenshots of all seven machines into one PDF file named \fn{screenshots.pdf} in
        this folder.
  \item Run \texttt{python3 test\_nfa.py} one last time to make sure everything is as you
        expect.
  \item Zip this entire folder and submit the zip as one file.
\end{enumerate}

\end{document}
